\SourceLeaf{137}
Statut (in der zweiten Interpretation) eigentlich \textit{unbedingt} übernommenen Verpflichtungen nicht mehr erfüllen kann.

Die vorstehenden Entwickelungen sollten die vitale Wichtigkeit der Progressionsnormirung bei einer Gesellschaft, deren Theilnehmerzahl sich bedeutend vergrössert, und damit die Wahrheit der S. [127] von mir aufgestellten Behauptung darthun. Um aber diesen Gegenstand von allen Seiten zu beleuchten, wird es nothwendig sein, dem Hergange der Entstehung jenes Statutsartikels Schritt vor Schritt zu folgen. Ich habe zu dem Zweck die betreffenden Acten sorgfältig gelesen, und wiederholt gelesen, und gebe daraus, soweit sie jenen Statutsartikel betreffen, einen Auszug. Ich werde dabei hin und wieder auch einige an sich untergeordnete Nebenumstände hervorzuheben haben, wenn sie etwas beitragen können, den Hergang bei diesen Verhandlungen begreiflicher zu machen. Im voraus will ich bemerken, dass 1794 der jährliche Beitrag 5 Thl. Gold betrug, die Witwenpension 110 Thl. Kassenmünze, und dass die vater- und mutterlosen Waisen die Pension bis zum vollendeten \(12^{\text{ten}}\) Jahre zu geniessen hatten.

In einem vom 30. Junius 1794 datirten an die Universität gerichteten Ministerial-Rescript, wurde unter Bezugnahme auf ein schon vor einiger Zeit von dem Könige der Witwenkasse gemachtes Geschenk von 1000 Thl.\footnote{Kassenmünze. Es war nach Ausweis der Rechnung für 1793 unter dem 18. April 1793 eingezahlt.}, die Anzeige von der Bewilligung eines zweiten Geschenks von 500 Thl. Gold gemacht, mit dem Beifügen, dass, wie bei diesen Geschenken die Absicht dahin gehe, zu einer baldigen Erhöhung der Pensionen hinzuwirken, gewärtigt werde, dass die Theilnehmer auch ihrerseits zur Erreichung dieses Zwecks beizutragen, und zu einer Erhöhung der jährlichen Beiträge von 5 Thl. Gold auf 10 Thl. Kassenmünze bereit sein würden. Nach den Berechnungen in einem anliegenden P. M. sei es nicht zweifelhaft, dass es füglich thunlich sei, schon jetzt eine Erhöhung der Pensionen in dem Maasse eintreten zu lassen, dass die sechs ältesten, anstatt der bisherigen 110 Thl., künftig 150 Thl. und alle übrigen jede 130 Thl. erhielten. Am Schlusse erbot man sich, falls die Kapitalien der Witwenkasse nicht alle vollkommen sicher placirt seien, die sichere Unterbringung zu 3 Procent bei öffentlichen oder städtischen Kassen zu veranlassen. Ein ähnliches Anerbieten war schon einmal, bei den Monitis zu der Jahresrechnung für 1792 gemacht worden.

Die beigefügte Anlage, deren Verfasser nicht genannt ist, im Detail durchzugehen, ist für meinen Zweck nicht nöthig. Aber ein paar Nebenumstände will ich herausheben.

I. Die Kapitalien der Witwenkssse, heisst es, seien zu ungleichem Zinsfuss ausgeliehen, einige\footnote{zu damaliger Zeit beinahe der dritte Theil des Kapitalvermögens.} zu 3 proc., andere höher. Weil aber der Zinsfuss leicht von allen Kapitalien auf 3 proc. heruntergehen könnte, und man bei zu machenden Ueberschlägen auf möglichst sichere Summen rechnen müsse, so wolle man bei den Rechnungen auch nicht mehr als 3 proc. voraussetzen.

Hieraus und aus dem eben angeführten Schlusse des Rescripts erklärt es sich, warum auch in den Verhandlungen bei der Universität für allen künftigen Kapitalzuwachs (ohne weitere Bemerkung) nur auf 3 proc. gerechnet ist; bloss für die schon vorhandenen und schon belegten Kapitale sind die Zinsen zu \(3\tfrac12\) proc. ausgeworfen. Vergl. hiemit die Anmerkung zu S. [131.]

II. Um einen Ueberschlag zu machen, auf welche Zahl von Witwen die Rechnung gestellt werden müsse, fährt der Verf. fort:

‘Will man nun nach den gemachten Erfahrungen annehmen, dass 3 stehende Ehen eine Witwe zu er-\par
\noindent ‘nähren haben, so würden die 26 verehelichten Professoren (unter der Gesammtzahl von 36) etwan 9 Witwen\par
\noindent ‘zu erhalten haben. Da es aber mehrere Gewissheit gewährt, wenn man den äussersten Fall zur Basis nimmt,\par
\noindent ‘so setze man lieber, dass gegen \(2\tfrac12\) Ehen \textit{eine} Witwe in Anschlag zu bringen, so dass also die bestehenden\par
\noindent ‘26 Professor-Ehen zu erhalten haben würden — 10 Witwen.’
